\documentclass[11pt]{article}

\usepackage[T1]{fontenc}
\usepackage{lmodern}
\usepackage[margin=1in]{geometry}
\usepackage{amsmath,amssymb,bm,mathtools}
\usepackage{booktabs,tabularx}
\usepackage{xcolor}
\usepackage{hyperref}

\setlength{\emergencystretch}{2em}
\allowdisplaybreaks

\hypersetup{
  colorlinks=true,
  linkcolor=blue!55!black,
  citecolor=blue!55!black,
  urlcolor=blue!55!black,
  pdftitle={A Working Review of h-Recursion in Two-Dimensional CFT},
  pdfauthor={Private working note}
}

\newcommand{\Vir}{\mathrm{Vir}}
\newcommand{\sVir}{\mathrm{sVir}}
\newcommand{\NS}{\mathrm{NS}}
\newcommand{\R}{\mathrm{R}}
\newcommand{\F}{\mathcal F}
\newcommand{\Hh}{\mathcal H}
\newcommand{\Vv}{\mathcal V}
\newcommand{\Res}{\operatorname*{Res}}
\newcommand{\half}{\tfrac12}
\newcommand{\bvec}{\bm}
\newcommand{\order}{\mathcal O}
\newcommand{\id}{\mathbf 1}
\newcommand{\ketbra}[2]{|#1\rangle\langle #2|}
\newcolumntype{L}[1]{>{\raggedright\arraybackslash}p{#1}}

\title{A Working Review of \(h\)-Recursion in Two-Dimensional CFT\\[0.4em]
\large Virasoro, Neveu--Schwarz, and Ramond blocks}
\author{Private working note---not for circulation}
\date{\today}

\begin{document}

\maketitle

\begin{abstract}
This note reviews the practical recipe for Zamolodchikov recursion in an
internal conformal weight.  The emphasis is on the data that a numerical
implementation actually needs: the Kac pole, inverse null norm, factorized
three-point matrix element, shifted module, elliptic or plumbing
normalization, and regular large-weight part.  We begin with Virasoro sphere
and torus blocks, then summarize the \(\mathcal N=1\) NS and Ramond
generalizations.  The Ramond discussion makes explicit how the \(G_0\)
ground-state doublet is hidden in the scalar recursions of the literature.
The final sections isolate the obstruction to applying the same method
directly to a genus-two theta graph.
\end{abstract}

\tableofcontents

\section{\texorpdfstring{The question answered by \(h\)-recursion}
{The question answered by h-recursion}}

Fix a chiral algebra, a pants decomposition, the central charge \(c\), all
external weights, and all internal weights except one, denoted \(h\).
A conformal-block coefficient obtained by descendant sewing is a rational
function of \(h\).  Its singularities occur when the internal Verma module
contains a null vector.  The block can therefore be reconstructed from
\begin{equation}
 \boxed{
 \text{regular part at large \(h\)}
 \quad+\quad
 \text{sum over Kac poles and their residues}.}
 \label{eq:one-line-summary}
\end{equation}
This is the meaning of \(h\)-recursion.  The residue is local to the two
vertices adjacent to the singular internal edge.  The regular part is not
local in this sense: it depends on the topology, channel, and normalization.
This distinction will be central at genus two.

At fixed order in the plumbing parameters, only finitely many Kac labels
contribute.  Once the regular part is known, the recursion is triangular and
efficient; it never requires construction or inversion of a large Gram
matrix.

\subsection{Three versions that should not be conflated}

The same meromorphic idea appears in three commonly used forms:
\begin{enumerate}
 \item A recursion for coefficients of a local coordinate expansion, such
 as the sphere \(z\)-series.
 \item An elliptic recursion after extracting a universal prefactor and
 replacing \(z\) by a nome \(q(z)\).  This is usually the fastest numerical
 representation of a four-point block.
 \item A plumbing recursion on a multipoint or higher-genus graph, with one
 plumbing parameter \(q_e\) for every internal edge.
\end{enumerate}
The poles and null-vector residues are the same representation-theory data.
The prefactor and regular part differ.

\section{Universal residue mechanism}
\label{sec:universal-residue}

Let \(\mathcal V_{c,h}\) be a Verma module with a basis
\(\{|h;I\rangle\}\) and Gram matrix
\begin{equation}
 B^{(n)}_{IJ}(c,h)=\langle h;I|h;J\rangle ,
 \qquad |I|=|J|=n .
\end{equation}
A sewn block coefficient has the schematic form
\begin{equation}
 F_n(c,h)
 =
 \sum_{|I|=|J|=n}
 \rho_L(I)\,
 \big(B^{(n)}(c,h)^{-1}\big)^{IJ}\,
 \rho_R(J).
 \label{eq:gram-definition}
\end{equation}
Suppose that at \(h=h_{r,s}(c)\) a normalized singular vector
\begin{equation}
 |\chi_{r,s}\rangle
 =
 D_{r,s}|h_{r,s}\rangle
\end{equation}
appears at level \(\ell_{r,s}\).  Descendants of this vector form a Verma
submodule with highest weight
\begin{equation}
 h'_{r,s}=h_{r,s}+\ell_{r,s}.
\end{equation}
Near the zero of the Kac determinant, the inverse Gram matrix has a
rank-one singular part.  Define the inverse null-norm slope by
\begin{equation}
 A_{r,s}(c)
 =
 \left[
 \lim_{h\to h_{r,s}}
 \frac{\langle\chi_{r,s}(h)|\chi_{r,s}(h)\rangle}
      {h-h_{r,s}}
 \right]^{-1}.
 \label{eq:inverse-null-slope}
\end{equation}
The Ward identities factorize the two adjacent three-point tensors:
\begin{align}
 \rho_L(D_{r,s}u)
 &=P^{(L)}_{r,s}\,\rho'_L(u),&
 \rho_R(D_{r,s}v)
 &=P^{(R)}_{r,s}\,\rho'_R(v).
 \label{eq:null-factorization}
\end{align}
The \(P_{r,s}\) are fusion polynomials, or finite matrices when the ground
space is not one-dimensional.  Equations
\eqref{eq:inverse-null-slope}--\eqref{eq:null-factorization} imply
\begin{equation}
 \Res_{h=h_{r,s}}\F
 =
 q^{\ell_{r,s}}\,
 A_{r,s}\,
 P^{(L)}_{r,s}P^{(R)}_{r,s}\,
 \F(h_{r,s}+\ell_{r,s}).
 \label{eq:universal-residue}
\end{equation}
This is the reusable kernel of every \(h\)-recursion discussed below.

\subsection{What must be fixed before quoting a residue}

A residue coefficient is meaningless until the following choices are
locked:
\begin{enumerate}
 \item the normalization of the singular operator \(D_{r,s}\);
 \item whether \(A_{r,s}\) is a slope with respect to \(h\), a momentum, or
 another parameter;
 \item the ordering and normalization of the two three-point forms;
 \item the local coordinate or elliptic prefactor;
 \item the branch convention for momenta such as \(b\), \(P\), or \(\beta\).
\end{enumerate}
A large fraction of apparent sign and factor-of-two disagreements in the
literature come from changing one item in this list.

\section{What is the regular boundary data?}
\label{sec:regular-boundary}

\subsection{Coefficient-level definition}

At a fixed plumbing level \(N\), suppose the reduced block coefficient is a
meromorphic function of one recursion variable \(x\):
\begin{equation}
 H_N(x)
 =
 U_N(x)
 +
 \sum_p\frac{R_{p,N}}{x-x_p}.
 \label{eq:coefficient-partial-fraction}
\end{equation}
The pole sum is fixed by null-vector factorization.  The function \(U_N\) is
the regular boundary data.  If \(H_N(x)\) is bounded as \(x\to\infty\), then
Liouville's theorem gives
\begin{equation}
 U_N=\lim_{x\to\infty}H_N(x).
 \label{eq:bounded-regular-part}
\end{equation}
If \(H_N\) grows as a polynomial of degree \(k\), then \(U_N(x)\) is a
polynomial of degree \(k\), and \(k+1\) asymptotic coefficients are needed.
In practice one chooses a prefactor so that the reduced block is bounded and
the regular part is as simple as possible.

At the level of the complete plumbing series,
\begin{equation}
 \Hh_\infty(\bvec q)
 =
 \sum_N \bvec q^{\,N} U_N
 \label{eq:complete-seed}
\end{equation}
is called the large-weight seed.  Thus ``the seed is one'' means
\begin{equation}
 U_0=1,\qquad U_{N>0}=0
\end{equation}
for the reduced block, not for the original unstripped conformal block.

\subsection{Why a prefactor is part of the boundary condition}

The original block may contain a known exponential or character-like
large-weight dependence:
\begin{equation}
 \F_x(\bvec q)
 =
 \mathcal P_x(\bvec q)\,\Hh_x(\bvec q).
 \label{eq:prefactor-and-reduced-block}
\end{equation}
One first derives \(\mathcal P_x\), then applies the partial-fraction
argument to \(\Hh_x\).  Changing \(\mathcal P_x\) changes the regular part.
Consequently the statement
\begin{equation}
 \lim_{x\to\infty}\Hh_x=1
 \end{equation}
is meaningful only together with the chosen elliptic or plumbing
normalization.

\subsection{How the large-weight limit is derived}
\label{sec:large-h-gram}

There are two standard derivations.

\paragraph{Semiclassical monodromy and uniformization.}
For the sphere four-point block, Zamolodchikov studied a classical limit in
which \(c\), the internal weight, and the external weights are large with
their ratios controlled.  The monodromy problem determines the leading
large-internal-weight dependence in terms of the elliptic nome.  The answer
is precisely the prefactor in
Eq.~\eqref{eq:vir-elliptic-prefactor}; after it is removed, the reduced
elliptic block tends to one.  The pillow representation gives a geometric
interpretation of the same factor.

\paragraph{Large-weight power counting in the descendant basis.}
At fixed \(c\),
\begin{equation}
 [L_n,L_{-n}]
 =
 2nL_0+\frac{c}{12}(n^3-n)
 \sim 2nh
 \label{eq:large-h-commutator}
\end{equation}
on a state of large weight \(h\).  For a PBW state \(L_{-A}|h\rangle\), let
\([A]\) be the number of Virasoro generators.  After Gram--Schmidt
orthogonalization,
\begin{equation}
 \|\ell_{-A}|h\rangle\|^2\sim h^{[A]}.
 \label{eq:large-h-norm}
\end{equation}
Off-diagonal inverse-Gram elements are suppressed.  If a three-point tensor
contains two large-weight legs and one fixed-weight primary, Ward identities
allow the descendant generators to pass through the fixed primary at
subleading cost.  Hence
\begin{equation}
 \frac{
 \rho(\ell_{-A}h,d,\ell_{-B}h)}
 {\|\ell_{-A}|h\rangle\|\,
  \|\ell_{-B}|h\rangle\|}
 \longrightarrow \delta_{AB}.
 \label{eq:large-h-diagonal-vertex}
\end{equation}
The boundary seed is then obtained by summing the descendant states that
survive this diagonal limit.

\subsection{The standard examples}

\paragraph{Sphere four-point block.}
The original \(z\)-block has a nontrivial large-\(h\) behavior.  After the
elliptic prefactor is extracted, one has
\begin{equation}
 \Hh_h(q)\longrightarrow1.
 \label{eq:sphere-h-boundary}
\end{equation}
The coefficientwise regular data are therefore
\(U_0=1\), \(U_{N>0}=0\).

\paragraph{Torus one-point block.}
In the trace, the large-\(h\) normalized matrix element of a fixed primary
approaches the descendant norm.  Every descendant therefore contributes
one, and the unstripped boundary is the nondegenerate Verma character:
\begin{equation}
 \F_h^d(q)
 \longrightarrow
 q^{h-c/24}
 \prod_{n=1}^{\infty}\frac1{1-q^n}.
 \label{eq:torus-character-boundary}
\end{equation}
After this character is divided out, the reduced seed is again one.

\paragraph{Torus \(N\)-point necklace channel.}
The regular part in one weight \(h_i\) is complicated.  Rather than compute
it directly, Cho--Collier--Yin set
\begin{equation}
 h_i=H+a_i,\qquad a_1=0,
 \label{eq:necklace-common-weight}
\end{equation}
and take \(H\to\infty\) with every difference \(a_i=h_i-h_1\) fixed.  Each
vertex then has two large legs and one fixed external primary, so
Eq.~\eqref{eq:large-h-diagonal-vertex} applies around the entire necklace.
All descendant partitions must agree around the cycle, giving
\begin{equation}
 \F_{\mathrm{necklace}}
 \longrightarrow
 \sum_A(q_1q_2\cdots q_N)^{|A|}
 =
 \prod_{n=1}^{\infty}
 \frac1{1-(q_1q_2\cdots q_N)^n}.
 \label{eq:necklace-boundary}
\end{equation}
The block is regarded as a meromorphic function of the single common
variable \(H\).  A Kac pole on edge \(i\) occurs at
\begin{equation}
 H=h_{r,s}-a_i.
\end{equation}
Thus all edge poles and one simultaneous boundary value determine a closed
one-variable recursion.  This avoids the need to know the separate regular
functions \(U_i(\{h_{j\ne i}\})\).

\paragraph{Sphere \(N\)-point linear channel.}
A similar one-parameter limit sends all internal weights and the two endpoint
external weights to infinity with their differences fixed.  Only the
level-zero term survives, so
\begin{equation}
 \F_{\mathrm{linear}}\longrightarrow1.
 \label{eq:linear-boundary}
\end{equation}
The price is that the recursion shifts some external weights as well as the
internal weights \cite{ChoCollierYin2017Recursive}.

\subsection{Superconformal large-weight boundaries}

For NS and R modules the same Gram--Schmidt argument includes fermionic
PBW modes.  In a channel where the large-weight vertex becomes diagonal,
the surviving-state sum is a super-Virasoro character:
\begin{equation}
 \sum_{A,K}Q^{|A|+|K|}
 =
 \prod_{n\ge1}\frac1{1-Q^n}
 \prod_{r\in\mathcal I_s}(1+\epsilon Q^r).
 \label{eq:super-character-boundary}
\end{equation}
Here \(\mathcal I_s\) is the appropriate NS or R mode lattice and
\(\epsilon\) records a possible \((-1)^F\) insertion.

In the Ramond case the power counting is naturally performed at
\(|\beta|\to\infty\).  The special parity basis makes the leading Gram
matrix diagonal.  Unstarred components approach a character, while a
component with an effective \(G_0\) insertion has an additional leading
factor proportional to \(\beta\).  Removing both the character and this
\(\beta\) factor produces the seed used in
Eq.~\eqref{eq:r-h-recursion}
\cite{HadaszJaskolskiSuchanek2012Toric}.

\subsection{Why the theta graph is different}

The successful limits above arrange that every trivalent tensor contains
\emph{two} large legs and \emph{one fixed} primary.  This makes the leading
tensor a two-point norm and leads to
Eq.~\eqref{eq:large-h-diagonal-vertex}.

On a genus-two theta graph, a simultaneous large-weight limit makes all
three legs of each trinion large.  The leading three-point tensor is then a
genuine trivalent oscillator vertex rather than a diagonal two-point norm.
The Kac residues do not determine this vertex.  To obtain a closed
\(h\)-recursion one must derive and resum the resulting oscillator network.
For an \(NRR\) vertex the network also retains the Ramond ground fibre and
the \(G_0\) sewing matrices.

This gives a concrete criterion for a new \(h\)-recursion:
\begin{equation}
 \boxed{
 \begin{gathered}
 \text{find a one-parameter correlated large-weight limit,}\\
 \text{prove that its trivalent large-weight tensor has a closed form,}\\
 \text{and factor it so that the reduced block has a bounded seed.}
 \end{gathered}}
 \label{eq:new-h-recursion-criterion}
\end{equation}
If no such closed trivalent seed is available, the local \(h\)-pole formulas
remain valid but do not form a complete algorithm.

\section{Virasoro four-point blocks on the sphere}
\label{sec:vir-sphere}

Use
\begin{equation}
 c=1+6Q^2,\qquad Q=b+b^{-1},
\qquad
 h(\lambda)=\frac{Q^2-\lambda^2}{4}.
 \label{eq:vir-parameterization}
\end{equation}
The degenerate weights are
\begin{equation}
 h_{r,s}
 =
 \frac{Q^2-(rb+s b^{-1})^2}{4},
 \qquad r,s\in\mathbb Z_{>0},
 \label{eq:vir-kac-weight}
\end{equation}
and the singular vector is at level
\(\ell_{r,s}=rs\).

For the channel
\(\langle\phi_4(\infty)\phi_3(1)\phi_2(z)\phi_1(0)\rangle\),
normalize
\begin{equation}
 \F_h(z)
 =
 z^{h-h_1-h_2}
 \left(1+\sum_{n\ge1}z^nF_n(h)\right).
 \label{eq:vir-local-block}
\end{equation}
The first coefficient,
\begin{equation}
 F_1(h)
 =
 \frac{(h+h_2-h_1)(h+h_3-h_4)}{2h},
 \label{eq:vir-level-one}
\end{equation}
already displays the mechanism: a pole at the level-one degenerate weight
\(h=0\), a factor from each adjacent three-point function, and a shifted
module.

\subsection{Coefficient recursion}

At each level,
\begin{equation}
 F_n(h)
 =
 U_n
 \sum_{\substack{r,s\ge1\\rs\le n}}
 \frac{R_{r,s}}{h-h_{r,s}}\,
 F_{n-rs}(h_{r,s}+rs),
 \label{eq:vir-coefficient-recursion}
\end{equation}
where
\begin{equation}
 R_{r,s}
 =
 A_{r,s}\,
 P_{r,s}(h_4,h_3)\,
 P_{r,s}(h_2,h_1).
 \label{eq:vir-residue-coefficient}
\end{equation}
The \(h\)-independent term \(U_n\) is fixed by the large-\(h\) asymptotic.
In a local \(z\)-series it is nontrivial; the elliptic normalization packages
the complete sequence \(U_n\) into a universal prefactor.

\subsection{Elliptic normalization}

Introduce
\begin{equation}
 q(z)=\exp\!\left[-\pi\frac{K(1-z)}{K(z)}\right].
 \label{eq:elliptic-nome}
\end{equation}
Zamolodchikov's normalization writes
\begin{align}
 \F_h(z)
 &=
 (16q)^{h-(c-1)/24}
 z^{(c-1)/24-h_1-h_2}
 (1-z)^{(c-1)/24-h_2-h_3}
 \notag\\[-0.2em]
 &\qquad\times
 \theta_3(q)^{(c-1)/2-4\sum_i h_i}
 \Hh_h(q).
 \label{eq:vir-elliptic-prefactor}
\end{align}
The reduced block has the recursion
\begin{equation}
 \boxed{
 \Hh_h(q)
 =
 1+
 \sum_{r,s\ge1}
 \frac{(16q)^{rs}R_{r,s}}
      {h-h_{r,s}}\,
 \Hh_{h_{r,s}+rs}(q).}
 \label{eq:vir-elliptic-recursion}
\end{equation}
The essential achievement of the elliptic prefactor is
\begin{equation}
 \lim_{h\to\infty}\Hh_h(q)=1.
 \label{eq:vir-large-h-seed}
\end{equation}
The recursion is efficient because a requested coefficient of \(q^N\)
receives contributions only from \(rs\le N\).

\subsection{Practical implementation}

To compute through \(q^N\):
\begin{enumerate}
 \item enumerate \((r,s)\) with \(rs\le N\);
 \item evaluate \(h_{r,s}\), \(A_{r,s}\), and the two fusion polynomials;
 \item recursively evaluate the shifted block only through order \(N-rs\);
 \item memoize by shifted weight and remaining order;
 \item multiply by the prefactor \eqref{eq:vir-elliptic-prefactor}.
\end{enumerate}
The recursion tree grows much more slowly than the descendant basis.

\section{Virasoro torus one-point and multipoint blocks}
\label{sec:vir-torus}

For a torus one-point block,
\begin{equation}
 \F_h^{d}(q)
 =
 q^{h-c/24}
 \sum_{n\ge0}q^n
 \sum_{|I|=|J|=n}
 \rho(h;I,d,h;J)\,
 (B_h^{(n)})^{IJ}.
 \label{eq:torus-one-point-definition}
\end{equation}
The same Kac pole appears on the traced module.  Both sides of the singular
state meet the same vertex, so the residue again contains two null-vector
three-point factors:
\begin{equation}
 \Res_{h=h_{r,s}}\F_h^d
 =
 q^{rs}A_{r,s}
 P_{r,s}^{(1)}(d,h'_{r,s})
 P_{r,s}^{(2)}(d,h_{r,s})
 \F_{h'_{r,s}}^d.
 \label{eq:torus-residue}
\end{equation}
After extracting the torus large-\(h\) character, the reduced block has seed
one and satisfies the same triangular pattern as
\eqref{eq:vir-elliptic-recursion}; see
Ref.~\cite{HadaszJaskolskiSuchanek2009Torus}.

\subsection{Linear and necklace channels}

For a sphere multipoint block in a linear channel or a torus multipoint
block in a necklace channel, every internal edge has two well-defined
adjacent vertices.  Applying \eqref{eq:universal-residue} edge by edge gives
\begin{equation}
 \Res_{h_e=h_{r,s}}
 \F_{\Gamma}
 =
 q_e^{rs}A_{r,s}
 P_{r,s}^{(v_L,e)}P_{r,s}^{(v_R,e)}
 \F_{\Gamma}(h_e\to h_{r,s}+rs).
 \label{eq:graph-local-h-residue}
\end{equation}
For these channels the simultaneous large-weight limit can be organized
recursively, and closed \(h\)-recursions are available
\cite{ChoCollierYin2017Recursive}.

\subsection{Why a local residue does not automatically give a closed recursion}

For several internal weights, meromorphy in one variable gives
\begin{equation}
 \F_\Gamma(\bvec h)
 =
 U_e(\{h_{e'}:e'\ne e\})
 +
 \sum_{r,s}
 \frac{\mathcal R_{r,s}^{(e)}}
      {h_e-h_{r,s}}\,
 \F_\Gamma(h_e\to h_{r,s}+rs).
 \label{eq:multivariable-partial-fraction}
\end{equation}
The residue is known, but the regular function \(U_e\) need not be.
At a theta vertex all three legs are internal.  Sending one weight to
infinity changes a three-point tensor whose other two weights remain
dynamical.  The large-\(h_e\) limit is therefore not the same universal
external-leg asymptotic that closes the sphere linear and torus necklace
recursions.  This is the basic higher-genus obstruction.

Cho--Collier--Yin solve the arbitrary Virasoro plumbing problem instead by
recursing in \(c\).  Its large-\(c\) regular part factorizes into a global
\(\mathfrak{sl}(2)\) block and a universal vacuum contribution.  This is why
the existence of graph-local \(h\)-pole residues should not be confused with
the existence of a complete genus-two \(h\)-recursion.

\section{\texorpdfstring{\(\mathcal N=1\) Neveu--Schwarz recursion}
{N=1 Neveu--Schwarz recursion}}
\label{sec:ns-recursion}

Use the ordinary super-Virasoro central charge
\begin{equation}
 c=\frac32+3Q^2,\qquad Q=b+b^{-1},
\qquad h(\lambda)=\frac{Q^2-\lambda^2}{8}.
 \label{eq:ns-parameterization}
\end{equation}
An NS singular vector is labelled by
\begin{equation}
 h^{\NS}_{r,s}
 =
 \frac{Q^2-(rb+s b^{-1})^2}{8},
 \qquad r+s\in2\mathbb Z,
 \qquad
 \ell_{r,s}=\frac{rs}{2}.
 \label{eq:ns-kac-weight}
\end{equation}
Because levels may be half-integral, an NS block has even and odd
superconformal components.  A null vector of odd parity maps an even block
to an odd block and conversely.

\subsection{Component form}

Let \(\pi\in\{\mathrm e,\mathrm o\}\) label the parity of the sewn
three-point structure.  The elliptic recursion takes the schematic form
\begin{equation}
 \Hh_h^\pi(q)
 =
 g^\pi(q)
 +
 \sum_{\substack{r,s\ge1\\r+s\ {\rm even}}}
 \frac{q^{rs/2}
 A^{\NS}_{r,s}
 P_{r,s}^{(L,\pi)}
 P_{r,s}^{(R,\pi)}}
 {h-h^{\NS}_{r,s}}\,
 \Hh_{h^{\NS}_{r,s}+rs/2}^{\,\pi+\mathrm{parity}(rs/2)}(q).
 \label{eq:ns-schematic-recursion}
\end{equation}
The precise powers of \(16\), signs, and seed functions \(g^\pi\) depend on
the chosen four-point component and elliptic prefactor
\cite{HadaszJaskolskiSuchanek2007NSBlocks,Suchanek2011EllipticRecursion}.

The conceptual changes relative to Virasoro recursion are only:
\begin{enumerate}
 \item the Kac level is \(rs/2\);
 \item NS fusion polynomials have even/odd versions;
 \item the recursion can mix the two component blocks;
 \item the large-\(h\) prefactor contains bosonic and fermionic oscillator
 characters.
\end{enumerate}

\subsection{Toric NS blocks}

The same residue mechanism applies to torus one-point blocks.  The trace can
include either \(1\) or \((-1)^F\), producing the NS and
\(\widetilde{\NS}\) spin structures.  The two blocks differ by signs for
half-integral levels, while the Kac data are unchanged
\cite{HadaszJaskolskiSuchanek2012Toric}.

\section{Ramond recursion: where the ground-state mixing goes}
\label{sec:r-recursion}

\subsection{Long Ramond ground space}

Parameterize a generic Ramond weight by
\begin{equation}
 h_\R(\beta)=\frac{c}{24}-\beta^2.
 \label{eq:r-weight}
\end{equation}
In a parity basis \((w_\beta^+,w_\beta^-)\),
\begin{equation}
 G_0w_\beta^\pm
 =
 i e^{\mp i\pi/4}\beta\,w_\beta^\mp,
 \qquad
 G_0^2=-\beta^2=h_\R-\frac{c}{24}.
 \label{eq:r-ground-action}
\end{equation}
After rephasing \(w_\beta^-\), this is equivalently
\begin{equation}
 G_0
 =
 \beta
 \begin{pmatrix}0&-1\\1&0\end{pmatrix},
 \qquad
 (-1)^F=
 \begin{pmatrix}1&0\\0&-1\end{pmatrix}.
 \label{eq:r-simple-matrices}
\end{equation}
At \(\beta=0\) the long module shortens.  One should perform the recursion at
generic \(\beta\) and construct the irreducible short quotient separately.

\subsection{Ramond Kac poles}

The positive-parity singular state occurs at
\begin{equation}
 \beta_{r,s}
 =
 \frac{rb+s b^{-1}}{2\sqrt2},
 \qquad r+s\in2\mathbb Z+1,
 \qquad \ell_{r,s}=\frac{rs}{2}.
 \label{eq:r-kac-momentum}
\end{equation}
Its null submodule is another long Ramond module with
\begin{equation}
 \beta'_{r,s}
 =
 \frac{(-1)^s}{2\sqrt2}
 \left(rb-sb^{-1}\right),
 \qquad
 h_\R(\beta'_{r,s})
 =
 h_\R(\beta_{r,s})+\frac{rs}{2}.
 \label{eq:r-shifted-momentum}
\end{equation}
The odd null partner is generated by \(G_0\):
\begin{equation}
 \chi^-_{r,s}
 =
 \frac{e^{i\pi/4}}{i\beta'_{r,s}}\,
 G_0\chi^+_{r,s}.
 \label{eq:r-null-partner}
\end{equation}

\subsection{\texorpdfstring{Why the published \(h\)-recursion is usually
scalar}{Why the published h-recursion is usually scalar}}

The toric construction of
Ref.~\cite{HadaszJaskolskiSuchanek2012Toric} starts with \(w_\beta^+\) but
allows \(G_0\) in the descendant strings.  It separates descendants into
even and odd fermion parity and defines
\begin{equation}
 F_{\beta,\mathrm e}^{\lambda(\pm)},\qquad
 F_{\beta,\mathrm o}^{\lambda(\pm)}.
\end{equation}
Here \(\mathrm e/\mathrm o\) is descendant parity, while \((\pm)\) labels
the two independent chiral R--R--NS three-point forms.  Ward identities give
\begin{equation}
 F_{\beta,\mathrm e}^{\lambda(\pm)}
 =
 \pm F_{\beta,\mathrm o}^{\lambda(\pm)},
 \label{eq:r-even-odd-relation}
\end{equation}
so it is enough to recurse two scalar blocks.  The \(2\times2\) ground
mixing has not disappeared: it has been diagonalized into these component
relations.

\subsection{\texorpdfstring{The Ramond \(h\)-recursion}
{The Ramond h-recursion}}

The natural meromorphic variable is \(\beta\), with poles at
\(\beta=\pm\beta_{r,s}\).  Combining the two poles gives a denominator
\begin{equation}
 h_\R(\beta)-h^\R_{r,s}
 =
 -(\beta-\beta_{r,s})(\beta+\beta_{r,s}).
 \label{eq:r-paired-poles}
\end{equation}
After extracting the appropriate large-\(\beta\) prefactor, the toric
recursion is
\begin{equation}
 \boxed{
 H_{\beta,\mathrm e}^{\lambda(\pm),f}
 =
 \delta_{f0}
 +
 \sum_{\substack{1<rs\le2f\\r+s\ {\rm odd}}}
 \frac{
 A^\R_{r,s}
 P^{r,s}_{c}\!\left[{}^{\Delta_\lambda}_{\pm\beta'_{r,s}}\right]
 P^{r,s}_{c}\!\left[{}^{\Delta_\lambda}_{\pm\beta_{r,s}}\right]}
 {h_\R(\beta)-h^\R_{r,s}}\,
 H_{\beta'_{r,s},\mathrm e}^{\lambda(\pm),f-rs/2}.}
 \label{eq:r-h-recursion}
\end{equation}
The phases from \(G_0\), the odd null partner, and the two ground states are
encoded in the two fusion polynomials and the \((\pm)\) label.

\subsection{\texorpdfstring{Blocks with an effective \(G_0\) insertion}
{Blocks with an effective G0 insertion}}

An odd Ramond plumbing modulus gives
\begin{equation}
 e^{\alpha G_0}=\id+\alpha G_0,
 \qquad
 \int d\alpha\,e^{\alpha G_0}=G_0.
 \label{eq:r-odd-sewing}
\end{equation}
When the \(G_0\) insertion can be moved to an adjacent vertex by a Ward
identity, it becomes a starred or odd chiral component.  Its leading
large-\(\beta\) behavior contains an explicit factor of \(\beta\).  In the
HJS toric conventions one finds schematically
\begin{equation}
 \F_{\beta}^{*\,\lambda(-)}
 =
 e^{i\pi/4}\sqrt2\,\beta\,
 q^{h_\R-c/24}
 \chi_\R(q)\,
 \Hh_{\beta}^{\lambda(-)}(q).
 \label{eq:r-g0-prefactor}
\end{equation}
After the \(\beta\) factor is extracted, the same reduced elliptic block
obeys the scalar recursion \eqref{eq:r-h-recursion}.  This is how known
low-topology \(h\)-recursions deal with the \(G_0\) matrix element.

\subsection{External Ramond fields versus an internal Ramond edge}

These two situations must be distinguished:
\begin{enumerate}
 \item A four-R sphere block in the \(R\times R\) channel propagates an NS
 module.  Its recursion uses NS Kac poles together with R--R--NS fusion
 polynomials.
 \item A block with an internal R module has the paired \(\beta\) poles,
 shifted momentum \(\beta'_{r,s}\), and ground doublet described above.
\end{enumerate}
The first problem tests Ramond external Ward identities but not the
Ramond propagator.

\section{\texorpdfstring{A practical \(h\)-recursion checklist}
{A practical h-recursion checklist}}
\label{sec:checklist}

For any proposed block, proceed in the following order.

\subsection*{Step 1: specify the block}
\begin{itemize}
 \item Choose the pants decomposition and local coordinates.
 \item Label every internal edge by Virasoro, NS, or R.
 \item Fix the ordered three-point form at every vertex.
 \item For R edges, fix parity basis, orientation, and spin lift.
\end{itemize}

\subsection*{Step 2: identify the meromorphic variable}
\begin{itemize}
 \item Virasoro or NS: normally the internal weight \(h\).
 \item Ramond: preferably \(\beta\), pairing the poles
 \(\pm\beta_{r,s}\).
\end{itemize}

\subsection*{Step 3: derive the pole kernel}
\begin{equation}
 \text{kernel}
 =
 \text{inverse null slope}
 \times
 \text{left fusion factor}
 \times
 \text{right fusion factor}.
\end{equation}
For an open R ground fibre, the last two factors are matrices and their
order matters.

\subsection*{Step 4: determine the shifted block}
\begin{center}
\begin{tabular}{@{}lll@{}}
\toprule
Sector & Kac level & Shift\\
\midrule
Virasoro & \(rs\) & \(h\to h_{r,s}+rs\)\\
NS & \(rs/2\) & \(h\to h^\NS_{r,s}+rs/2\)\\
R & \(rs/2\) & \(\beta\to\beta'_{r,s}\)\\
\bottomrule
\end{tabular}
\end{center}

\subsection*{Step 5: derive the regular part}

This is mandatory.  A list of poles is not yet a recursion.  One must
derive, rather than guess,
\begin{equation}
 \F_{\infty}
 =
 \lim_{\text{internal weight}\to\infty}
 \left(\text{properly reduced block}\right).
 \label{eq:regular-part-requirement}
\end{equation}
For standard sphere and torus elliptic blocks this limit is one.  For
multipoint and higher-genus blocks it may be a nontrivial function or
finite-dimensional matrix.

\subsection*{Step 6: truncate by plumbing degree}

At total plumbing degree \(N\), retain only poles whose null level does not
exceed the remaining degree.  Memoization should use the full shifted
weight data, component label, spin lift, and remaining multidegree.

\subsection*{Step 7: perform collision-safe evaluation}

At rational or self-dual values of \(b\), distinct labels can give the same
pole.  Terms must first be combined algebraically as Laurent or
partial-fraction objects.  Evaluating each resonant term separately and
then adding generally loses the finite part.

\section{Common failure modes}
\label{sec:failure-modes}

\begin{center}
\begin{tabularx}{\textwidth}{@{}L{0.27\textwidth}X@{}}
\toprule
Failure & Diagnosis\\
\midrule
Poles without a seed &
The result is only a residue identity, not a closed recursion.\\

Using \(z\) where a formula assumes \(q\) &
Powers of \(16q\), theta functions, and the large-weight seed become
incorrect.\\

Confusing \(A^{(h)}_{r,s}\) and \(A^{(\beta)}_{r,s}\) &
The two differ by the Jacobian
\(-1/(2\beta_{r,s})\).\\

Discarding the NS parity label &
Odd null vectors map even and odd component blocks into one another.\\

Treating the R ground doublet as two unrelated modules &
\(G_0\) relates the states and fixes relative phases.\\

Projecting the R ground fibre too early &
Spin-lift, \(G_0\), and GSO matrices may not commute; orientation signs are
lost.\\

Substituting \(\beta=0\) into a long-module formula &
The Ramond module shortens and requires an irreducible quotient.\\

Calling a chiral block a correlator &
The nonchiral pairing, structure constants, momentum integrals, spin sum,
and string measure remain to be included.\\
\bottomrule
\end{tabularx}
\end{center}

\section{Genus two: what the review tells us}
\label{sec:genus-two}

\subsection{The theta graph}

A genus-two vacuum surface can be obtained by gluing two trinions along
three edges.  In an \(\mathcal N=1\) theory the chiral sector assignments are
\begin{equation}
 NNN
 \qquad\text{or}\qquad
 NRR
\end{equation}
at each vertex.  For an \(NRR\) theta graph, denote the weights by
\begin{equation}
 h_N,\qquad \beta_1,\qquad\beta_2
\end{equation}
and the plumbing parameters by
\(q_N,q_{R,1},q_{R,2}\).

The edge-local \(h\)-pole data are known:
\begin{align}
 \Res_{h_N=h^\NS_{r,s}}\F_\Theta
 &=
 q_N^{rs/2}\,
 \mathcal R^\NS_{r,s}\,
 \F_\Theta(h_N\to h^\NS_{r,s}+rs/2),
 \notag\\
 \Res_{\beta_a=\pm\beta_{r,s}}\F_\Theta
 &=
 q_{R,a}^{rs/2}\,
 \mathbf P^{(a)}_{L,r,s}\,
 \F_\Theta(\beta_a\to\beta'_{r,s})\,
 \mathbf P^{(a)}_{R,r,s}.
 \label{eq:nrr-local-residues}
\end{align}
The Ramond formula is matrix-valued until the two odd sewings are
contracted with
\begin{equation}
 G_0^{(1)}\otimes G_0^{(2)}.
 \label{eq:two-g0}
\end{equation}

\subsection{The missing ingredient}

The unknown object is the simultaneous regular part,
\begin{equation}
 \mathbf U_\Theta
 =
 \text{regular part of }
 \F_\Theta(h_N,\beta_1,\beta_2)
 \text{ in the internal weights}.
 \label{eq:theta-regular-part}
\end{equation}
The scalar torus relations
\(F_{\mathrm e}=\pm F_{\mathrm o}\) cannot simply be applied independently
to both R edges, because the two R lines meet at the same ordered R--R--NS
tensor.  One must first derive the joint Ward identities and the resulting
finite matrix of component blocks.

Consequently,
\begin{equation}
 \boxed{
 \text{known Ramond \(h\)-pole kernel}
 \;\not\Rightarrow\;
 \text{known genus-two Ramond \(h\)-recursion}.}
 \label{eq:not-yet-recursion}
\end{equation}

\subsection{Two plausible routes}

\paragraph{Route A: derive a genuine multi-weight \(h\)-seed.}
Choose a basis of R--R--NS component tensors, extract the explicit
\(\beta_1\beta_2\) from the two \(G_0\) sewings, and derive the joint
large-\((h_N,\beta_1,\beta_2)\) asymptotic.  If the result resums to a known
oscillator determinant, the local residues
\eqref{eq:nrr-local-residues} close a matrix-valued \(h\)-recursion.

\paragraph{Route B: use genus-two \(c\)-recursion.}
Keep \(h_N,\beta_1,\beta_2\) fixed while varying \(c\), so that
\(G_0^2=-\beta_a^2\) remains finite.  The Kac residues are again edge-local.
The remaining problem is a large-\(c\) \(NRR\) seed, which should follow
from the contracted super-Virasoro oscillator algebra.  This route is the
closer analogue of the arbitrary-genus Virasoro construction of
Ref.~\cite{ChoCollierYin2017Recursive}.

\subsection{\texorpdfstring{Fixed-\(\beta\) Ramond \(c\)-poles}{Fixed-beta Ramond c-poles}}
\label{sec:fixed-beta-c-poles}

The fixed-\(\beta\) formulation deserves to be separated from a
fixed-\(h_R\) \(c\)-recursion.  Write
\begin{equation}
 h_R(c,\beta)=\frac{c}{24}-\beta^2
 \end{equation}
and vary \(c\) with \(\beta\) held fixed.  Then
\begin{equation}
 G_0^2=-\beta^2,
 \qquad
 q^{h_R-c/24}=q^{-\beta^2}
 \label{eq:fixed-beta-benefits}
\end{equation}
are independent of \(c\).  In particular, the level-zero denominator
\((h_R-c/24)^{-1}\) is simply \(-\beta^{-2}\); it does not produce an
additional pole at \(c=24h_R\).

A positive-level Ramond pole occurs when
\begin{equation}
 \beta=\epsilon\beta_{r,s}(b),
 \qquad
 \beta_{r,s}(b)=\frac{rb+s b^{-1}}{2\sqrt2},
 \qquad
 \epsilon=\pm1,
 \qquad
 r+s\ {\rm odd}.
 \label{eq:fixed-beta-degeneracy}
\end{equation}
Equivalently, the pole value \(b_*\) solves
\begin{equation}
 r b_*^2-\epsilon\,2\sqrt2\,\beta b_*+s=0,
 \qquad
 c^{R}_{r,s}(\beta)
 =
 \frac32+3\left(b_*+b_*^{-1}\right)^2.
 \label{eq:fixed-beta-c-pole}
\end{equation}
The quadratic branches and the labels related by
\(b\leftrightarrow b^{-1}\), \(r\leftrightarrow s\) must be organized once
and continued consistently so that the same \(c\)-pole is not counted
twice.

The denominator in the Ramond \(h\)-recursion becomes, along the
fixed-\(\beta\) path,
\begin{equation}
 h_R(c,\beta)-h^R_{r,s}(c)
 =
 \beta_{r,s}(c)^2-\beta^2.
 \label{eq:fixed-beta-denominator}
\end{equation}

\subsubsection{\texorpdfstring{Residue recipe: first the null module, then the
\(c\)-plane}{Residue recipe: first the null module, then the c-plane}}

It is useful to separate two logically distinct operations.  The
representation-theoretic residue is first determined in a local
weight--momentum coordinate.  It is then pulled back to the fixed-\(\beta\)
\(c\)-plane.

At level \(f\), a block coefficient has the schematic Gram-matrix form
\begin{equation}
 F_f
 =
 \sum_{I,J}
 \rho_L(\nu_I)\,
 [B_f(c,\beta)^{-1}]^{IJ}\,
 \rho_R(\nu_J).
 \label{eq:fixed-beta-gram-coefficient}
\end{equation}
At \(\beta=\epsilon\beta_{r,s}\), the primitive Ramond null multiplet
\(\chi^\pm_{r,s}\) appears at level
\(\ell_{r,s}=rs/2\).  Only the inverse Gram matrix restricted to the
descendant submodule generated by this multiplet is singular.  Thus it is
more precise to say that the primitive null direction creates the pole,
while all descendants of the null multiplet resum into a block for the
shifted Ramond module.

Choose the normalization of the singular-vector operator \(D_{r,s}\) once
and define its inverse-norm coefficient by
\begin{equation}
 A^{(h),R}_{r,s}
 =
 \lim_{h\to h^R_{r,s}}
 (h-h^R_{r,s})\,
 \langle D_{r,s}w\mid D_{r,s}w\rangle^{-1}.
 \label{eq:fixed-beta-null-norm}
\end{equation}
The three-point matrix elements of \(D_{r,s}w\) factor into left and right
Ramond fusion matrices,
\begin{equation}
 \rho_L(D_{r,s}w)=\mathbf P_{L,r,s}\rho_L(w_{\beta'_{r,s}}),
 \qquad
 \rho_R(D_{r,s}w)=\rho_R(w_{\beta'_{r,s}})
 \mathbf P_{R,r,s}.
 \label{eq:fixed-beta-fusion-factorization}
\end{equation}
Keeping the two Ramond ground indices open automatically includes both
\(\chi^+_{r,s}\) and its \(G_0\)-related partner \(\chi^-_{r,s}\).

\paragraph{Published Ramond kernel.}
The data above are not new.  They were worked out for four-point blocks
with R intermediate states in Ref.~\cite{Suchanek2011EllipticRecursion}
and for toric one-point blocks in
Ref.~\cite{HadaszJaskolskiSuchanek2012Toric}.  In the normalization
\begin{equation}
 D_{r,s}=L_{-1}^{rs/2}+\cdots ,
\end{equation}
the weight-residue normalization is
\begin{equation}
 A^{(h),R}_{r,s}
 =
 2^{rs-2}\left(rb+\frac{s}{b}\right)
 \prod_{\substack{m=1-r,\ldots,r\\ n=1-s,\ldots,s\\
                   m+n\in2\mathbb Z\\(m,n)\neq(0,0)}}
 \left(mb+n b^{-1}\right)^{-1}.
 \label{eq:published-Ah-R}
\end{equation}
The coefficient used for a pole in \(\beta\) in
Ref.~\cite{Suchanek2011EllipticRecursion} already includes the first
Jacobian:
\begin{equation}
 A^{(\beta),R}_{r,s}
 =
 -\frac{A^{(h),R}_{r,s}}{2\beta_{r,s}}
 =
 -2^{rs-\frac32}
 \prod_{\substack{m=1-r,\ldots,r\\ n=1-s,\ldots,s\\
                   m+n\in2\mathbb Z\\(m,n)\neq(0,0)}}
 \left(mb+n b^{-1}\right)^{-1}.
 \label{eq:published-Abeta-R}
\end{equation}
Because \(r+s\) is odd, the factor \((m,n)=(r,s)\) is absent from the even
sublattice automatically, and the two expressions obey exactly
\(A^{(\beta),R}_{r,s}=-A^{(h),R}_{r,s}/(2\beta_{r,s})\).
The odd-sublattice condition printed below the corresponding product in
Ref.~\cite{Suchanek2011EllipticRecursion} is a typo: it fails already for
the exact level-one \(2\times2\) Ramond Gram matrix.  The even sublattice
reproduces the residues at both \(\beta_{2,1}\) and \(\beta_{1,2}\).

For an adjacent NS weight
\begin{equation}
 h_N=\frac{Q^2}{8}-\frac{\lambda^2}{8},
 \qquad Q=b+b^{-1},
\end{equation}
define
\begin{align}
 X^{r,s}_{\rm e}(x)
 &=
 \prod_{\substack{p=1-r+2k,\ q=1-s+2l\\
                   0\leq k\leq r-1,\ 0\leq l\leq s-1\\
                   k+l\ {\rm even}}}
 \frac{x-pb-qb^{-1}}{2\sqrt2},
 \notag\\
 X^{r,s}_{\rm o}(x)
 &=
 \prod_{\substack{p=1-r+2k,\ q=1-s+2l\\
                   0\leq k\leq r-1,\ 0\leq l\leq s-1\\
                   k+l\ {\rm odd}}}
 \frac{x-pb-qb^{-1}}{2\sqrt2},
 \label{eq:published-Xeo-R}\\
 P^{r,s}_{c,\sigma}(\beta;\lambda)
 &=
 X^{r,s}_{\rm e}(2\sqrt2\beta-\sigma\lambda)\,
 X^{r,s}_{\rm o}(2\sqrt2\beta+\sigma\lambda),
 \qquad \sigma=\pm1.
 \label{eq:published-fusion-R}
\end{align}
These are the two Ramond fusion-polynomial components denoted
\(P^{rs}_c[{}^{\Delta_\lambda}_{\pm\beta}]\) in the toric paper.

In the component basis of the four-point paper, the residue at
\(\beta=\epsilon\beta_{r,s}\) takes the universal form
\begin{equation}
 \Res_{\beta=\epsilon\beta_{r,s}}
 \mathbf F
 =
 \epsilon\,A^{(\beta),R}_{r,s}\,
 \mathbf P^{(\epsilon)}_{L,r,s}\,
 \mathbf F(\beta'_{r,s})\,
 \mathbf P^{(\epsilon)}_{R,r,s},
 \qquad \epsilon=\pm1.
 \label{eq:published-beta-kernel}
\end{equation}
Here each \(\mathbf P^{(\epsilon)}\) is selected from
\eqref{eq:published-fusion-R} with the component and orientation signs
specified by the corresponding R--R--NS three-point form.  Equation
\eqref{eq:published-beta-kernel}, including these sign changes, is the
local kernel to import for every Ramond plumbing edge.

At fixed \(c\), the same denominator may be resolved into the two momentum
poles:
\begin{equation}
 \frac{1}{h_R(\beta)-h^R_{r,s}}
 =
 -\frac{1}{(\beta-\beta_{r,s})(\beta+\beta_{r,s})},
 \qquad
 \Res_{\beta=\epsilon\beta_{r,s}}
 =
 -\frac{1}{2\epsilon\beta_{r,s}}.
 \label{eq:fixed-beta-beta-residue}
\end{equation}
In the \(c\)-recursion, however, \(\beta\) itself is not varied.  The local
coordinate at the pole is
\begin{equation}
 u_{r,s}(c)=\epsilon\beta_{r,s}(c)-\beta,
 \qquad
 u_{r,s}(c_*)=0,
 \label{eq:fixed-beta-local-coordinate}
\end{equation}
and
\begin{equation}
 \beta_{r,s}(c)^2-\beta^2
 =
 2\beta\,u_{r,s}(c)+O(u_{r,s}^2).
 \label{eq:fixed-beta-pullback}
\end{equation}
Consequently, the conversion from the momentum residue to the \(c\)-plane
is entirely the Jacobian of this change of local coordinate.

Therefore the pole-conversion Jacobian is
\begin{align}
 J^R_{r,s}(\beta)
 &=
 \left[
 \frac{\partial\,\beta_{r,s}(c)^2}{\partial c}
 \right]^{-1}_{c=c^R_{r,s}(\beta)}
 \notag\\
 &=
 \left.
 \frac{dc/db}{d\beta_{r,s}^2/db}
 \right|_{b=b_*}
 =
 \left.
 \frac{
 24(b+b^{-1})(1-b^{-2})}
 {(rb+s b^{-1})(r-s b^{-2})}
 \right|_{b=b_*}.
 \label{eq:fixed-beta-jacobian}
\end{align}
Equivalently, if one imports the published \(\beta\)-residue coefficient
\eqref{eq:published-Abeta-R} directly, then
\begin{equation}
 J^R_{r,s}A^{(h),R}_{r,s}
 =
 -\left.
 \frac{A^{(\beta),R}_{r,s}}
 {\partial_c\beta_{r,s}}
 \right|_{c=c^R_{r,s}(\beta)}.
 \label{eq:published-kernel-pullback}
\end{equation}
This identity is the complete modification required to turn the known
Ramond \(h\)-recursion pole kernel into a fixed-\(\beta\) \(c\)-plane
kernel; it does not require a new singular-vector calculation.

With ground indices kept open, the corresponding edge residue is
\begin{equation}
 \begin{aligned}
 \Res_{c=c^R_{r,s}(\beta_e)}
 \mathbf F_\Gamma
 ={}&
 q_e^{rs/2}\,
 J^R_{r,s}(\beta_e)\,
 A^{(h),R}_{r,s}\,
 \mathbf P^{(e)}_{L,r,s}\,
 \mathbf F_\Gamma
 \left(
 c^R_{r,s}(\beta_e),
 \beta_e\to\beta'_{r,s}
 \right)
 \mathbf P^{(e)}_{R,r,s}.
 \end{aligned}
 \label{eq:fixed-beta-c-residue}
\end{equation}
The shifted momentum \(\beta'_{r,s}\) is
Eq.~\eqref{eq:r-shifted-momentum}, evaluated at \(b=b_*\).

For an \(NRR\) theta graph the natural analytic variables are therefore
\begin{equation}
 \left(c;h_N,\beta_1,\beta_2\right),
 \end{equation}
with \(h_N,\beta_1,\beta_2\) fixed in the \(c\)-recursion.  A pole on one R
edge shifts only its \(\beta\); the other R momentum remains fixed even
though its absolute conformal weight changes with the recursively evaluated
central charge.

The remaining boundary datum is now the single limit
\begin{equation}
 \mathbf F_{\Theta,\infty}^{NRR}
 =
 \lim_{\substack{c\to\infty\\h_N,\beta_1,\beta_2\ {\rm fixed}}}
 \mathbf F_\Theta(c;h_N,\beta_1,\beta_2).
 \label{eq:fixed-beta-large-c-seed}
\end{equation}
This is simpler analytically than a multivariate large-\(h\) boundary
problem, but it is not yet known: the R weights scale as \(c/24\), so the
seed is a hybrid large-\(c\) limit rather than the ordinary light
superconformal block.

\subsection{Sphere \(R\)-channel certificate}
\label{sec:sphere-r-c-certificate}

Before using the local kernel on a higher-genus graph, consider the mixed
sphere block
\begin{equation}
 \left\langle
 V^{\NS}_4(\infty)\,
 V^{\R}_3(1)\,
 V^{\R}_2(z)\,
 V^{\NS}_1(0)
 \right\rangle ,
 \qquad
 \widehat H_{\beta}^{\eta_3,\eta_2}(c;q),
 \label{eq:sphere-r-test-block}
\end{equation}
with an R intermediate module.  The analytic variables are
\begin{equation}
 \left(c;h_1,\beta_2,\beta_3,h_4;\beta,\eta_3,\eta_2\right);
 \label{eq:sphere-r-fixed-data}
\end{equation}
all variables after the semicolon are held fixed while \(c\) moves.

\paragraph{Canonical pole table.}
It is enough to retain the unordered labels
\(r>s\geq1\), \(r+s\) odd, and work on the positive sheet
\(\beta=\beta_{r,s}(b)\).  Its two roots are
\begin{equation}
 b_{r,s}^{(\alpha)}
 =
 \frac{\sqrt2\,\beta+
 \alpha\sqrt{2\beta^2-rs}}{r},
 \qquad
 c_{r,s}^{(\alpha)}
 =
 \frac32+3\left(
 b_{r,s}^{(\alpha)}+
 \frac{1}{b_{r,s}^{(\alpha)}}\right)^2,
 \qquad \alpha=\pm1 .
 \label{eq:sphere-r-two-c-poles}
\end{equation}
The description \(\beta=-\beta_{r,s}(b)\) gives the same two \(c\)-poles
after \(b\mapsto-b\).  Likewise, \(r\leftrightarrow s\) is accompanied by
\(b\leftrightarrow b^{-1}\).  Thus summing either of these equivalent
descriptions would double count the poles.  On the canonical sheet no
endpoint-sign flip occurs; if the negative sheet is used instead, both
\(\eta_2\) and \(\eta_3\) flip as in the published \(\beta\)-recursion.

At a stored root \(b_*\),
\begin{equation}
 \partial_c\beta_{r,s}
 =
 \left.
 \frac{r-sb^{-2}}
 {12\sqrt2\,(b+b^{-1})(1-b^{-2})}
 \right|_{b=b_*}.
 \label{eq:sphere-r-db-dc}
\end{equation}
Consequently the coefficient residue is
\begin{equation}
 \begin{split}
 \mathcal R_{r,s}^{(\alpha;\eta_3,\eta_2)}
 ={}&
 -\frac{16^{rs/2}A^{(\beta),R}_{r,s}(b_*)}
 {\partial_c\beta_{r,s}(b_*)}
 P_{r,s}^{\eta_3}(\beta_3,h_4;b_*)\,
 P_{r,s}^{\eta_2}(\beta_2,h_1;b_*).
 \end{split}
 \label{eq:sphere-r-c-residue}
\end{equation}
The shifted internal parameter is
\begin{equation}
 \beta'_{r,s}(b_*)
 =
 (-1)^s\frac{rb_*-s/b_*}{2\sqrt2}.
 \label{eq:sphere-r-beta-prime}
\end{equation}

\paragraph{The necessary regularization at \(c=\infty\).}
The raw HJS elliptic series is not bounded coefficient by coefficient at
fixed \(\beta\).  Already at level one its regular polynomial is
\begin{equation}
 [q]H_{\beta}^{\eta_3,\eta_2}\big|_{\rm reg}
 =
 \frac c3
{}+8\beta^2-8\beta_2^2-8\beta_3^2
-8h_1-8h_4-\frac12 .
 \label{eq:sphere-r-raw-level-one-seed}
\end{equation}
The leading powers through the first computed orders are generated by the
Ramond oscillator factor
\begin{equation}
 \mathcal E_R(c,q)
 =
 \left[
 \prod_{n=1}^{\infty}
 \frac{1+q^n}{1-q^n}
 \right]^{c/6}.
 \label{eq:sphere-r-oscillator-factor}
\end{equation}
We therefore define
\(H=\mathcal E_R\,\widehat H\).  Since \(\mathcal E_R\) is entire in \(c\)
and independent of the internal momentum and component labels, the pole
kernel is unchanged.  Its first three regular coefficients are exactly
\begin{equation}
 S_0=1,\qquad
 S_1(\beta)
 =
 8\beta^2-8\beta_2^2-8\beta_3^2
 -8h_1-8h_4-\frac12 .
 \label{eq:sphere-r-normalized-seed-one}
\end{equation}
If
\begin{equation}
 X=\beta^2-\beta_2^2-\beta_3^2-h_1-h_4,
\end{equation}
the next coefficient obtained directly from the fixed-\(\beta\)
large-\(c\) limit is
\begin{equation}
 S_2(\beta)
 =
 32X^2-16\beta^2+12\beta_2^2+12\beta_3^2
 -4h_1-4h_4-\frac38 .
 \label{eq:sphere-r-normalized-seed-two}
\end{equation}
Both \(S_1\) and \(S_2\) are independent of
\((\eta_3,\eta_2)\).
At arbitrary order the normalized recursion is
\begin{equation}
 \begin{split}
 [q^n]\widehat H_\beta(c)
 ={}&S_n(\beta)
 +\sum_{\substack{r>s\geq1,\ r+s\ {\rm odd}\\rs/2\leq n}}
 \sum_{\alpha=\pm1}
 \frac{\mathcal R_{r,s}^{(\alpha)}}
 {c-c_{r,s}^{(\alpha)}}\,
 [q^{n-rs/2}]
 \widehat H_{\beta'_{r,s}}
 \left(c_{r,s}^{(\alpha)}\right).
 \end{split}
 \label{eq:sphere-r-fixed-beta-c-recursion}
\end{equation}
The endpoint component labels in
\eqref{eq:sphere-r-fixed-beta-c-recursion} are unchanged on the canonical
positive sheet.

\paragraph{Level-one certificate.}
In the basis
\((L_{-1}w^+,G_{-1}G_0w^+)\), the exact Gram matrix is
\begin{equation}
 B_1=
 \begin{pmatrix}
 2h&-\frac32\beta^2\\
 -\frac32\beta^2&-\beta^2(2h+c/4)
 \end{pmatrix},
 \qquad h=\frac c{24}-\beta^2.
 \label{eq:sphere-r-level-one-gram}
\end{equation}
Its nontrivial Kac polynomial is
\begin{equation}
 K_1(c,\beta)
 =
 c^2-30\beta^2c+144\beta^4+81\beta^2.
 \label{eq:sphere-r-level-one-kac}
\end{equation}
The two zeros of \(K_1\) are precisely
\(c_{2,1}^{(+)}\) and \(c_{2,1}^{(-)}\).
Contracting \(B_1^{-1}\) with the two Ward vectors, or instead using
\eqref{eq:sphere-r-fixed-beta-c-recursion} with the two residues
\eqref{eq:sphere-r-c-residue} and the seed
\eqref{eq:sphere-r-normalized-seed-one}, gives the same answer for all four
\((\eta_3,\eta_2)\).  The numerical residue test approaches the predicted
value with relative error below \(2\times10^{-9}\).
At \(q^2\), the same recursion includes the shifted-\(\beta'\) level-one
tail and the two \((4,1)\) poles.  Together with
\eqref{eq:sphere-r-normalized-seed-two}, it reproduces the full HJS
coefficient for all four component choices.

This completes the pole part of the sphere R-channel \(c\)-recursion.  The
remaining derivation problem is sharply isolated: determine the closed
hybrid large-\(c\) seed \(S_n(\beta)\) for \(n\geq3\).  Substituting an
ordinary scalar \(\mathfrak{osp}(1|2)\) global block here would be
unjustified because the R weights scale as \(c/24\).

\subsection{\texorpdfstring{The fixed-\(\beta\) contracted Ramond algebra}
{The fixed-beta contracted Ramond algebra}}
\label{sec:contracted-r-algebra}

The sphere calculation suggests that the regular block should itself be
regarded as a conformal block for the algebra which survives at
\(c=\infty\).  This ``smaller algebra'' is not literally a subalgebra
obtained by discarding most super-Virasoro generators.  It is an
In\"on\"u--Wigner contraction of the full Ramond algebra, together with a
renormalized limit of its three-point intertwiners.

\paragraph{Starting algebra and fixed-\(\beta\) grading.}
In the ordinary-\(c\) convention,
\begin{align}
 [L_m,L_n]
 &=(m-n)L_{m+n}
 +\frac{c}{12}(m^3-m)\delta_{m+n,0},
 \notag\\
 [L_m,G_n]
 &=\left(\frac m2-n\right)G_{m+n},
 \notag\\
 \{G_m,G_n\}
 &=2L_{m+n}
 +\frac c3\left(m^2-\frac14\right)\delta_{m+n,0},
 \qquad m,n\in\mathbb Z .
 \label{eq:r-supervir-algebra}
\end{align}
The correct finite zero-mode operator in the fixed-\(\beta\) limit is
\begin{equation}
 D=L_0-\frac c{24}.
 \label{eq:contracted-D}
\end{equation}
On the Ramond ground fiber
\(\mathcal R_\beta^{(0)}\),
\begin{equation}
 D|\beta,a\rangle=-\beta^2|\beta,a\rangle,
 \qquad
 G_0^2=D,
 \label{eq:contracted-ground}
\end{equation}
where \(a\) is the open two-dimensional Ramond ground index.  Thus both
\(D\) and \(G_0\) remain unscaled and finite as \(c\to\infty\).

\paragraph{Rescaled nonzero modes.}
For every positive integer \(n\), define
\begin{equation}
 \begin{aligned}
 a_n&=\sqrt{\frac{12}{cn^3}}\,L_n,
 &\qquad
 a_n^\dagger&=\sqrt{\frac{12}{cn^3}}\,L_{-n},
 \\
 \psi_n&=\sqrt{\frac{3}{cn^2}}\,G_n,
 &
 \psi_n^\dagger&=\sqrt{\frac{3}{cn^2}}\,G_{-n}.
 \end{aligned}
 \label{eq:contracted-oscillators}
\end{equation}
These are normalized versions of the original superconformal modes, not
new local fields.  Here \(\dagger\) denotes the creation-mode partner in
the chiral BPZ pairing; no Hilbert-space reality condition is being
imposed.  The \(c\to\infty\) statements below are understood
coefficientwise on fixed-level descendants.  On a state at finite level
\(\ell\), \(D=-\beta^2+\ell\) remains \(O(c^0)\); no claim is made about a
simultaneous limit in which the descendant level grows with \(c\).

Acting on the ground fiber, the normalized one-oscillator descendants are
exactly
\begin{align}
 a_n^\dagger|\beta,a\rangle
 &=
 \frac{L_{-n}}{\sqrt{cn^3/12}}|\beta,a\rangle,
 \notag\\
 \psi_n^\dagger|\beta,a\rangle
 &=
 \frac{G_{-n}}{\sqrt{cn^2/3}}|\beta,a\rangle .
 \label{eq:normalized-original-descendants}
\end{align}
\begin{samepage}
Indeed, their unnormalized diagonal norms are
\begin{align}
 \langle\beta,a|L_nL_{-n}|\beta,b\rangle
 &=
 \left(\frac{cn^3}{12}-2n\beta^2\right)
 \langle\beta,a|\beta,b\rangle,
 \notag\\
 \langle\beta,a|G_nG_{-n}|\beta,b\rangle
 &=
 \left(\frac{cn^2}{3}-2\beta^2\right)
 \langle\beta,a|\beta,b\rangle.
 \label{eq:original-mode-large-c-norms}
\end{align}
\end{samepage}
The factors in Eq.~\eqref{eq:contracted-oscillators} are therefore the
unique positive-mode normalizations, up to finite phases, which produce
unit leading Gram matrix.  At equal mode number the corresponding exact
diagonal brackets are
\begin{align}
 [a_n,a_n^\dagger]
 &=1+\frac{24D}{cn^2},
 \notag\\
 \{\psi_n,\psi_n^\dagger\}
 &=1+\frac{6D}{cn^2}.
 \label{eq:contracted-diagonal-finite-c}
\end{align}
On any fixed-level family of states, brackets at unequal mode number are
suppressed by \(c^{-1/2}\).  The \(c\to\infty\) oscillator algebra is
therefore
\begin{align}
 [a_n,a_m^\dagger]&=\delta_{nm},
 & [a_n,a_m]&=[a_n^\dagger,a_m^\dagger]=0,
 \notag\\
 \{\psi_n,\psi_m^\dagger\}&=\delta_{nm},
 & \{\psi_n,\psi_m\}
 =\{\psi_n^\dagger,\psi_m^\dagger\}=0,
 \label{eq:contracted-free-oscillators}
\end{align}
with all boson--fermion supercommutators zero.

The grading operator is not central.  Its surviving action is
\begin{align}
 [D,a_n]&=-n a_n,
 & [D,a_n^\dagger]&=n a_n^\dagger,
 \notag\\
 [D,\psi_n]&=-n\psi_n,
 & [D,\psi_n^\dagger]&=n\psi_n^\dagger.
 \label{eq:contracted-D-action}
\end{align}
Most importantly, the Ramond zero mode does not decouple:
\begin{align}
 [G_0,a_n]&=-\sqrt n\,\psi_n,
 &
 [G_0,a_n^\dagger]&=\sqrt n\,\psi_n^\dagger,
 \notag\\
 \{G_0,\psi_n\}&=\sqrt n\,a_n,
 &
 \{G_0,\psi_n^\dagger\}&=\sqrt n\,a_n^\dagger,
 \notag\\
 G_0^2&=D,
 & [D,G_0]&=0.
 \label{eq:contracted-G0-action}
\end{align}
Equations
\eqref{eq:contracted-free-oscillators}--\eqref{eq:contracted-G0-action}
are the complete contracted Ramond algebra, which we denote by
\begin{equation}
 \mathfrak A_{\R,\infty}
 =
 \left(
 \mathfrak h_{\rm bos}\oplus\mathfrak h_{\rm fer}
 \right)\rtimes\langle D,G_0\rangle .
 \label{eq:contracted-r-algebra-name}
\end{equation}
It is an infinite super-Heisenberg algebra acted on by a common Ramond
supercharge.  In particular, every nonzero mode survives after the
appropriate rescaling; the limit is simpler because their nonlinear
mode-changing interactions vanish, not because only finitely many modes
remain.

\paragraph{Fock realization of the contracted Ramond module.}
Let \(\Gamma_\beta\) be the odd operator on
\(\mathcal R_\beta^{(0)}\) satisfying
\begin{equation}
 \Gamma_\beta^2=-\beta^2.
 \label{eq:contracted-ground-clifford}
\end{equation}
The tensor product with the oscillator Fock space is graded, so
\(\Gamma_\beta\) anticommutes with the \(\psi_n,\psi_n^\dagger\) and
commutes with the \(a_n,a_n^\dagger\).  A representation of
\(\mathfrak A_{\R,\infty}\) is
\begin{align}
 N&=
 \sum_{n=1}^{\infty}n
 \left(a_n^\dagger a_n+\psi_n^\dagger\psi_n\right),
 \notag\\
 D&=-\beta^2+N,
 \notag\\
 G_0&=
 \Gamma_\beta+
 \sum_{n=1}^{\infty}\sqrt n\,
 \left(a_n^\dagger\psi_n+\psi_n^\dagger a_n\right).
 \label{eq:contracted-fock-realization}
\end{align}
The graded cross terms in \(G_0^2\) vanish, and each oscillator pair
contributes
\(n(a_n^\dagger a_n+\psi_n^\dagger\psi_n)\); hence
\(G_0^2=D\) holds on the entire contracted module.

A diagonal oscillator basis is
\begin{equation}
 |\boldsymbol k,\boldsymbol\varepsilon;\beta,a\rangle
 =
 \prod_{n\geq1}
 (a_n^\dagger)^{k_n}
 (\psi_n^\dagger)^{\varepsilon_n}
 |\beta,a\rangle,
 \qquad
 k_n\in\mathbb Z_{\geq0},\quad
 \varepsilon_n\in\{0,1\},
 \label{eq:contracted-fock-basis}
\end{equation}
with oscillator level
\begin{equation}
 |\boldsymbol k,\boldsymbol\varepsilon|
 =
 \sum_{n\geq1}n(k_n+\varepsilon_n).
 \label{eq:contracted-level}
\end{equation}
Apart from the fixed ground-fiber metric, its Gram matrix is diagonal with
bosonic norm
\(\prod_n k_n!\).  This is the first reason the contracted block can be
evaluated without constructing super-Virasoro Gram matrices.

\begin{samepage}
\paragraph{The conformal block of the contracted algebra.}
Let
\(\rho_{L,\infty}\) and \(\rho_{R,\infty}\) denote the two renormalized
N--R--R three-point intertwiners adjacent to the internal R edge.  Keeping
the Ramond ground indices open, define
\begin{align}
 \mathbf S_\beta(q)
 ={}&
 \sum_{\boldsymbol k,\boldsymbol\varepsilon}
 q^{|\boldsymbol k,\boldsymbol\varepsilon|}
 \rho_{L,\infty}
 \left(
 |\boldsymbol k,\boldsymbol\varepsilon;\beta\rangle
 \right)
 \left(B_{\infty}^{-1}\right)_
 {\boldsymbol k,\boldsymbol\varepsilon}
 \rho_{R,\infty}
 \left(
 |\boldsymbol k,\boldsymbol\varepsilon;\beta\rangle
 \right)
 \notag\\
 ={}&
 \left\langle\mathcal V_{L,\infty}\right|
 q^N
 \left|\mathcal V_{R,\infty}\right\rangle .
 \label{eq:contracted-r-block}
\end{align}
\end{samepage}
The primary propagation \(q^{-\beta^2}\) has already been extracted.
Equation \eqref{eq:contracted-r-block} is the precise meaning of the
``conformal block for the smaller algebra.''  For the sphere channel,
\begin{equation}
 \mathbf S_\beta(q)
 =
 \lim_{\substack{c\to\infty\\
 h_1,h_4,\beta,\beta_2,\beta_3\ {\rm fixed}}}
 \widehat{\mathbf H}_\beta(c;q)
 =
 \sum_{n=0}^{\infty}\mathbf S_n(\beta)q^n .
 \label{eq:contracted-r-block-limit}
\end{equation}
The scalar component computed in
Eqs.~\eqref{eq:sphere-r-normalized-seed-one} and
\eqref{eq:sphere-r-normalized-seed-two} is independent of the endpoint
component labels through \(q^2\).

\paragraph{Renormalized intertwiners are additional data.}
The contracted commutation relations alone do not determine
\(\rho_{L,\infty}\) and \(\rho_{R,\infty}\).  They must be obtained by
putting the exact N--R--R Ward identities in the variables
\eqref{eq:contracted-oscillators}, holding all \(\beta\)'s and the adjacent
NS weight fixed, and then taking \(c\to\infty\).  The leading coherent
sources grow as \(\sqrt c\); their Gaussian sewing produces the universal
factor \(\mathcal E_R(c,q)\) in
\eqref{eq:sphere-r-oscillator-factor}.  Thus the canonical statement is
about the sewn functional,
\begin{equation}
 \left\langle\mathcal V_{L,\infty}\right|q^N
 \left|\mathcal V_{R,\infty}\right\rangle
 :=
 \lim_{c\to\infty}
 \mathcal E_R(c,q)^{-1}
 H_\beta(c;q).
 \label{eq:contracted-renormalized-vertices}
\end{equation}
Splitting the counterterm into a separate square root assigned to each
vertex is a convention, whereas the combined sewn limit above is
unambiguous.

Since different positive mode numbers form independent free
super-oscillators, the contracted Ward identities are expected to reduce
the block to a product of small mode kernels,
\begin{equation}
 \mathbf S_\beta(q)
 =
 \left\langle v_L\right|
 \overrightarrow{\prod_{n=1}^{\infty}}
 \mathbf K_n(q^n)
 \left|v_R\right\rangle .
 \label{eq:contracted-mode-product}
\end{equation}
The arrow records the ground-fiber and fermionic ordering.  Once
\(\mathbf K_n\) is derived, truncating
\eqref{eq:contracted-mode-product} at total \(q\)-degree \(N\) gives the
first \(N\) seed coefficients exactly; it is not a truncation of the
original super-Virasoro Gram problem.

\paragraph{Derivation target.}
The next representation-theoretic task is therefore finite and concrete:
\begin{enumerate}
 \item contract the two ordered N--R--R Ward identities using
       Eqs.~\eqref{eq:contracted-oscillators} and
       \eqref{eq:contracted-G0-action};

 \item isolate the coherent \(O(c)\) overlap producing
       \(\mathcal E_R\);

 \item determine the finite ground-fiber kernels \(\mathbf K_n\);

 \item check that their product reproduces \(S_1\) and \(S_2\);

 \item use the resulting exact product as the regular input at every call
       of the fixed-\(\beta\) \(c\)-recursion.
\end{enumerate}
This approach addresses the accuracy problem directly: the regular block
is derived as an exact contracted-algebra block and is not estimated by
ordinary descendant-level truncation.

\subsection{\texorpdfstring{Complete fixed-\(\beta\) \(c\)-recursion
recipe}{Complete fixed-beta c-recursion recipe}}
\label{sec:complete-fixed-beta-recipe}

The following data are sufficient to turn the pole identities into an
algorithm.  We state them for the \(NRR\) theta block, whose natural reduced
chiral block is
\begin{equation}
 \mathbf F_\Theta
 =
 \mathbf F_\Theta
 \left(c;h_N,\beta_1,\beta_2;q_N,q_1,q_2\right).
 \label{eq:nrr-natural-variables}
\end{equation}
The reduction must first extract the primary propagation factors
\begin{equation}
 q_N^{h_N-c/24},
 \qquad q_1^{-\beta_1^2},
 \qquad q_2^{-\beta_2^2},
 \label{eq:nrr-primary-prefactors}
\end{equation}
and fix the normalization of the two ordered R--R--NS tensors.  All R
ground indices are kept open until the final sewing or GSO contraction.

\paragraph{1. Meromorphy and the subtraction degree.}
At each fixed plumbing multidegree \(\boldsymbol n\), one must prove that the
reduced coefficient is meromorphic in \(c\), with only Kac poles, and
determine its growth at \(c=\infty\).  If it is bounded, then
\begin{equation}
 \mathbf F_{\Theta,\boldsymbol n}(c)
 =
 \mathbf U_{\Theta,\boldsymbol n}
 +
 \sum_{p\in\mathcal P_{\boldsymbol n}}
 \frac{\mathbf R_{p,\boldsymbol n}}{c-c_p}.
 \label{eq:nrr-partial-fraction}
\end{equation}
If it grows polynomially, \(\mathbf U_{\Theta,\boldsymbol n}\) must be
replaced by the complete subtraction polynomial.  This analytic statement,
not merely the list of Kac zeros, is what closes a \(c\)-recursion.

\paragraph{2. Pole tables.}
For the NS edge, hold \(h_N\) fixed and solve
\begin{equation}
 h_N=h^{\NS}_{r,s}(c^{\NS}_{r,s}(h_N)),
 \qquad r+s\ {\rm even}.
\end{equation}
Use
\begin{equation}
 J^{\NS}_{r,s}(h_N)
 =
 \left[-\partial_c h^{\NS}_{r,s}(c)\right]^{-1}
 _{c=c^{\NS}_{r,s}(h_N)}.
 \label{eq:nrr-ns-jacobian}
\end{equation}
For an R edge \(e\), hold \(\beta_e\) fixed and solve
\begin{equation}
 \beta_e=\epsilon\beta_{r,s}(b_*),
 \qquad
 c^*_{e,r,s,\epsilon}=c(b_*),
 \qquad r+s\ {\rm odd}.
 \label{eq:nrr-r-pole-table}
\end{equation}
Store the chosen \(b_*\) together with the \(c\)-pole: it is needed to
evaluate \(A_{r,s}\), the fusion polynomials, and \(\beta'_{r,s}\).
Equivalent \(b\leftrightarrow b^{-1}\) branches are canonicalized before
the pole is inserted.

\paragraph{3. Universal local residue kernels.}
For a pole on the NS edge,
\begin{equation}
 \begin{aligned}
 \Res_{c=c^{\NS}_{r,s}(h_N)}\mathbf F_\Theta
 ={}&
 q_N^{rs/2}
 J^{\NS}_{r,s}A^{\NS}_{r,s}
 \Sigma^{(L)}_{r,s}\Sigma^{(R)}_{r,s}\,
 \mathbf F_\Theta
 \left(c^{\NS}_{r,s};
 h_N+\frac{rs}{2},\beta_1,\beta_2\right).
 \end{aligned}
 \label{eq:nrr-complete-ns-residue}
\end{equation}
For a pole on R edge \(e\),
\begin{equation}
 \begin{aligned}
 \Res_{c=c^*_{e,r,s,\epsilon}}\mathbf F_\Theta
 ={}&
 q_e^{rs/2}
 J^{R}_{r,s}(\beta_e;\epsilon)A^{(h),R}_{r,s}
 \mathbf P^{(L,e;\epsilon)}_{r,s}
 \\
 &\times
 \mathbf F_\Theta
 \left(c^*_{e,r,s,\epsilon};
 \beta_e\longrightarrow\beta'_{r,s}(b_*)\right)
 \mathbf P^{(R,e;\epsilon)}_{r,s}.
 \end{aligned}
 \label{eq:nrr-complete-r-residue}
\end{equation}
All labels not displayed in the shifted argument are spectators.  One may
equivalently use the published \(\beta\)-residue coefficient and replace
\(J^R A^{(h),R}\) by
\(-A^{(\beta),R}/\partial_c\beta_{r,s}\).

\paragraph{4. Component and sewing ledger.}
Every recursive state must include
\begin{equation}
\begin{gathered}
 (\text{edge sector},\ \text{half-edge order},\ \pm\text{ three-form},
 \ \text{ground parity},\\
 \text{spin lift},\ \text{R orientation},\
 \text{remaining multidegree}).
\end{gathered}
 \label{eq:nrr-recursion-state}
\end{equation}
The published fusion polynomials determine the scalar entries, but the
ledger determines which entries and phases occur.  Odd Ramond plumbing
matrices and GSO projectors are applied only after the open-index block has
been assembled.

\paragraph{5. Regular boundary data.}
The one genuinely new representation-theory input is
\begin{equation}
 \mathbf U_\Theta(\boldsymbol q)
 =
 \lim_{\substack{c\to\infty\\h_N,\beta_1,\beta_2\ {\rm fixed}}}
 \mathbf F_\Theta(c;h_N,\beta_1,\beta_2;\boldsymbol q).
 \label{eq:nrr-complete-seed}
\end{equation}
It should be derived by contracting the large-\(c\) super-Virasoro
oscillator algebra while keeping \(G_0^2=-\beta_e^2\) finite.  A direct
Gram/Ward calculation through the first few multidegrees is needed only to
identify and certify this seed, not as the final evaluation method.

\paragraph{6. Triangular recursion and base cases.}
At coefficient \((n_N,n_1,n_2)\), include a pole on edge \(e\) only when
\(rs/2\leq n_e\), and call the shifted block at
\(n_e-rs/2\).  Set negative multidegrees to zero and fix the level-zero
matrix from the chosen R--R--NS three-point normalization.  Memoization keys
must contain \(c\), all shifted labels, component data, and the remaining
multidegree.

\paragraph{7. Collision-safe evaluation.}
Distinct Kac labels or quadratic branches can yield the same \(c_p\),
especially at \(b=1\).  Equal pole keys must be merged as exact
partial-fraction or Laurent objects before substituting the physical
central charge.  The recursion is performed at generic long
\(\beta_e\neq0\); the short quotient at \(\beta_e=0\) is a separate final
limit.

\paragraph{8. Certification.}
The implementation should pass, in increasing order of cost:
\begin{enumerate}
 \item direct Gram/Ward coefficients at the first two nonzero
 multidegrees;
 \item the published sphere and torus R-channel recursions after the
 appropriate graph degeneration;
 \item \(b\leftrightarrow b^{-1}\) and \(\beta\leftrightarrow-\beta\)
 component relations;
 \item equality of the two torus two-point plumbing channels;
 \item factorization and spin-structure checks after nonchiral sewing.
\end{enumerate}

Consequently, the currently known package is
\begin{equation}
 \boxed{
 \text{pole positions}+\text{Jacobians}+\text{null norms}
 +\text{fusion matrices}+\text{label shifts}.}
\end{equation}
The missing package is
\begin{equation}
 \boxed{
 \text{fixed-\(\beta\) large-\(c\) seed}
 +\text{fully explicit higher-genus component ledger}.}
\end{equation}

\subsection{Role of Virasoro-primary decomposition}

Restricting a super-Virasoro module to Virasoro modules gives an infinite
tower of Virasoro-primary families.  This may help derive the large-\(c\) or
large-weight seed by resumming known Virasoro genus-two blocks.  It is
unlikely to be an efficient final algorithm because it introduces an
infinite sum, nontrivial branching three-point tensors, and spurious
Virasoro poles that cancel only after complete supermultiplets are
recombined.

\section{Established results and open tasks}

\begin{center}
\begin{tabularx}{\textwidth}{@{}L{0.28\textwidth}L{0.24\textwidth}X@{}}
\toprule
Object & Status & Comment\\
\midrule
Virasoro sphere four-point \(h\)-recursion &
Established &
Elliptic seed equals one.\\

Virasoro torus one-point and necklace recursion &
Established &
Universal character prefactor and local Kac residues.\\

Arbitrary Virasoro plumbing graph &
Established via \(c\)-recursion &
Large-\(c\) seed is global block times vacuum factor.\\

\(\mathcal N=1\) NS sphere and torus \(h\)-recursion &
Established &
Requires even/odd component bookkeeping.\\

Ramond sphere and torus \(h\)-recursion &
Established in fixed components &
\(G_0\) mixing is encoded in parity relations, fusion polynomials, and
explicit \(\beta\) prefactors.\\

Genus-two \(NNN\) theta block &
Not closed &
Local poles are known; the multivariate regular seed is not.\\

Genus-two \(NRR\) theta block &
Not closed &
Requires open R ground fibres, two \(G_0\) sewings, coupled component
recursion, and a regular seed.\\
\bottomrule
\end{tabularx}
\end{center}

\section{Recommended immediate derivation}

The shortest useful calculation is not a high-level descendant expansion.
It is the following finite problem:
\begin{enumerate}
 \item Choose an explicit ordered basis for the two independent R--R--NS
 chiral tensors.
 \item Write \(G_0\), \((-1)^F\), and the two spin-lift sewing maps as
 \(2\times2\) matrices.
 \item Derive the Ward-identity action of a null embedding on this tensor
 basis.
 \item Determine whether the two R-edge residue matrices can be
 simultaneously reduced to scalar sectors.
 \item Compute the joint large-\(\beta_1,\beta_2\) leading tensor and decide
 whether the explicit factor \(\beta_1\beta_2\) leaves a finite seed.
\end{enumerate}
If step 4 succeeds, the published scalar Ramond \(h\)-recursion generalizes
almost directly.  If it fails, the correct object is a small matrix-valued
recursion.  In either case, step 5 decides whether \(h\)-recursion closes or
whether genus-two \(c\)-recursion is necessary.

\appendix

\section{Basis-independent form of Ramond sewing}

Let \(\mathcal G_\beta\) be the \(1|1\) Ramond ground fibre, \(B_0\) its BPZ
metric, and \(S_\epsilon\) the spin-lift map.  The fixed-odd-modulus sewing
operator is
\begin{equation}
 \mathcal K_\R(q,\alpha;\epsilon)
 =
 q^{L_0-c/24}\,
 B^{-1}\,
 S_\epsilon\,
 e^{\alpha G_0}.
 \label{eq:basis-independent-r-sewing}
\end{equation}
After Berezin integration,
\begin{equation}
 \int d\alpha\,
 \mathcal K_\R(q,\alpha;\epsilon)
 =
 q^{L_0-c/24}\,
 B^{-1}\,
 S_\epsilon G_0.
 \label{eq:integrated-r-sewing}
\end{equation}
The order of \(S_\epsilon\), \(G_0\), and the BPZ map is part of the
definition.  A scalar component recursion is a projection of
\eqref{eq:basis-independent-r-sewing}, not a replacement for it.

\section{Notation dictionary}

\begin{center}
\begin{tabularx}{\textwidth}{@{}L{0.25\textwidth}L{0.25\textwidth}X@{}}
\toprule
Symbol & Meaning & Warning\\
\midrule
\(h\) or \(\Delta\) & Conformal weight & Both notations occur in the
literature.\\
\(q(z)\) & Elliptic nome & Not the same as a local sphere coordinate \(z\).\\
\(A_{r,s}\) & Inverse null-norm slope & Its value depends on the
differentiation variable.\\
\(P_{r,s}\) & Null-vector fusion polynomial & Ordering and parity conventions
matter.\\
\(\pi=\mathrm e,\mathrm o\) & Superconformal component parity & Not a spin
structure by itself.\\
\((\pm)\) in R blocks & Chiral R--R--NS tensor structure & Not identical to
the ground-state parity label.\\
\(\beta\) & HJS Ramond momentum & For physical super-Liouville momentum,
typically \(\beta=iP/\sqrt2\).\\
\(\beta'_{r,s}\) & Momentum of shifted null submodule & It replaces
\(\beta\) after taking a Ramond residue.\\
\bottomrule
\end{tabularx}
\end{center}

\begin{thebibliography}{99}

\bibitem{Zamolodchikov1987Recursion}
A.~B. Zamolodchikov,
\emph{Conformal symmetry in two-dimensional space: recursion representation
of conformal block},
Theor.\ Math.\ Phys.\ \textbf{73} (1987) 1088--1093.

\bibitem{HadaszJaskolskiSuchanek2009Torus}
L.~Hadasz, Z.~Jaskolski, and P.~Suchanek,
\emph{Recursive representation of the torus one-point conformal block},
JHEP \textbf{01} (2010) 063,
\href{https://arxiv.org/abs/0911.2353}{arXiv:0911.2353}.

\bibitem{ChoCollierYin2017Recursive}
M.~Cho, S.~Collier, and X.~Yin,
\emph{Recursive representations of arbitrary Virasoro conformal blocks},
JHEP \textbf{09} (2017) 053,
\href{https://arxiv.org/abs/1703.09805}{arXiv:1703.09805}.

\bibitem{HadaszJaskolskiSuchanek2007NSBlocks}
L.~Hadasz, Z.~Jaskolski, and P.~Suchanek,
\emph{Recursion representation of the Neveu--Schwarz superconformal block},
JHEP \textbf{03} (2007) 032,
\href{https://arxiv.org/abs/hep-th/0611266}{hep-th/0611266}.

\bibitem{HadaszJaskolskiSuchanek2008RamondBlocks}
L.~Hadasz, Z.~Jaskolski, and P.~Suchanek,
\emph{Elliptic recurrence representation of the \(N=1\) superconformal
blocks in the Ramond sector},
JHEP \textbf{11} (2008) 060,
\href{https://arxiv.org/abs/0810.1203}{arXiv:0810.1203}.

\bibitem{Suchanek2011EllipticRecursion}
P.~Suchanek,
\emph{Elliptic recursion for 4-point superconformal blocks and bootstrap in
\(N=1\) SLFT},
JHEP \textbf{02} (2011) 090,
\href{https://arxiv.org/abs/1012.2974}{arXiv:1012.2974}.

\bibitem{HadaszJaskolskiSuchanek2012Toric}
L.~Hadasz, Z.~Jaskolski, and P.~Suchanek,
\emph{Recurrence relations for toric \(N=1\) superconformal blocks},
JHEP \textbf{09} (2012) 122,
\href{https://arxiv.org/abs/1207.5740}{arXiv:1207.5740}.

\bibitem{BelavinGeiko2018CRecursion}
V.~A. Belavin and R.~V. Geiko,
\emph{\(c\)-recursion for multi-point superconformal blocks: NS sector},
JHEP \textbf{08} (2018) 112,
\href{https://arxiv.org/abs/1806.09563}{arXiv:1806.09563}.

\bibitem{Witten2012SuperRiemannSurfaces}
E.~Witten,
\emph{Notes on super Riemann surfaces and their moduli},
\href{https://arxiv.org/abs/1209.2459}{arXiv:1209.2459}.

\end{thebibliography}

\end{document}
